\documentclass[aps,pra,reprint,superscriptaddress,nofootinbib,longbibliography]{revtex4-2}

\usepackage{amsmath,amssymb}
\usepackage{bm}

\begin{document}

\title{Comment on ``Entropic uncertainty principle for mixed states'':\
An exact qutrit counterexample to the extended-MUB conjecture}

\author{Xiaojian Zeng}
\email{xiaojianzeng@foxmail.com}
\affiliation{PBC School of Finance, Tsinghua University, Beijing 100190, China}

\date{\today}

\begin{abstract}
Rotundo and Schwonnek conjectured that a condition involving only the second
singular value of a doubly stochastic measurement-overlap matrix makes the
uniform vector a global maximizer of a related induced norm.  We give an exact
counterexample within the physically realizable qutrit subset.  Two complete
rank-one projective measurements are related by a rational real orthogonal
matrix, so their overlap is orthostochastic and hence unistochastic.  At
$\lambda=3/4$ and $\mu=1/3$, the proposed singular-value condition holds with a
strict exact margin, but the positive vector $(6,10,5)^{\mathsf T}$ produces a
norm ratio strictly larger than the mutually unbiased-basis value.  The two
strict inequalities reduce to a quadratic-radical comparison and a positive
integer certificate.  Thus local stability at the uniform vector does not
imply the conjectured global formula, even for two legal qutrit measurements.
The proved induced-norm uncertainty relation itself is unaffected.
\end{abstract}

\maketitle

\section{Context and claim}

For two complete rank-one projective measurements $X$ and $Y$ in dimension
$d$, let
\begin{equation}
 C_{ij}=\operatorname{Tr}(X_iY_j)=|U_{ij}|^2,
 \qquad U\in U(d),
 \label{eq:overlap}
\end{equation}
where the columns of $U$ express the $Y$ basis in the $X$ basis.  Rotundo and
Schwonnek derived, at entropy weight $\alpha=1$, the uncertainty relation
\begin{equation}
 \lambda H_X(\rho)+\mu H_Y(\rho)
 \geq S(\rho)-\log\|C\|_{r\to s},
 \label{eq:proved-bound}
\end{equation}
where $r=1/\mu$ and $s=1/(1-\lambda)$, and proposed an analytic formula for
the norm in an ``extended MUB regime''
\cite{RotundoSchwonnek2024}.  In the nontrivial domain
$0<\lambda,\mu<1$ and $\lambda+\mu>1$, their conjecture states that
\begin{equation}
 \frac{(1-\mu)(1-\lambda)}{\mu\lambda}\geq \sigma_2(C)^2
 \quad\Longrightarrow\quad
 \|C\|_{r\to s}=d^{\,1-\lambda-\mu}.
 \label{eq:conjecture}
\end{equation}
Here $\sigma_2(C)$ is the second singular value.  Because $C$ is entrywise
nonnegative, its induced norm may be written as
\begin{equation}
 \|C\|_{r\to s}
 =\max_{x\in\mathbb R_+^d\setminus\{0\}}
 \frac{\|Cx\|_s}{\|x\|_r}.
 \label{eq:norm}
\end{equation}
The uniform vector always gives the right-hand side of
Eq.~\eqref{eq:conjecture} as a lower bound.  The supplemental analysis of
Ref.~\cite{RotundoSchwonnek2024} shows that the singular-value condition makes
this vector a local maximum; the missing step is global optimality.

We disprove Eq.~\eqref{eq:conjecture} already for $d=3$, while retaining the
physical restriction $C_{ij}=|U_{ij}|^2$.  No relaxation to an arbitrary
doubly stochastic matrix is used.

\section{Exact physical counterexample}

Take the real matrix
\begin{equation}
 U=\begin{pmatrix}
  1/3&-14/15&2/15\\
  -2/3&-1/3&-2/3\\
  2/3&2/15&-11/15
 \end{pmatrix}.
 \label{eq:U}
\end{equation}
Direct rational multiplication gives $U^{\mathsf T}U=I_3$ and
$\det U=1$.  Hence $U$ is an exact unitary matrix on $\mathbb C^3$, and its
columns define a complete rank-one projective measurement.  Equation
\eqref{eq:overlap} gives
\begin{equation}
 C=|U|^2=\begin{pmatrix}
  1/9&196/225&4/225\\
  4/9&1/9&4/9\\
  4/9&4/225&121/225
 \end{pmatrix}.
 \label{eq:C}
\end{equation}
Every row and column sums to one.  More importantly, Eq.~\eqref{eq:U}
certifies that $C$ is orthostochastic, and therefore unistochastic.

Choose
\begin{equation}
 \lambda=\frac34,\qquad \mu=\frac13,
 \qquad r=3,\qquad s=4.
 \label{eq:parameters}
\end{equation}
The parameter conditions are satisfied because
$\lambda+\mu=13/12>1$.  To check spectral admissibility exactly, the
characteristic polynomial of $C^{\mathsf T}C$ factors as
\begin{equation}
 \det(tI-C^{\mathsf T}C)
 =(t-1)\frac{50625t^2-33714t+49}{50625}.
 \label{eq:charpoly}
\end{equation}
Its nontrivial eigenvalues are
\begin{equation}
 \eta_{\pm}=\frac{1873\pm64\sqrt{849}}{5625},
 \qquad \sigma_2(C)^2=\eta_+.
 \label{eq:eigenvalues}
\end{equation}
For Eq.~\eqref{eq:parameters}, the threshold in
Eq.~\eqref{eq:conjecture} equals $2/3$, and
\begin{equation}
 \frac23-\sigma_2(C)^2
 =\frac{1877-64\sqrt{849}}{5625}>0.
 \label{eq:spectral-margin}
\end{equation}
The last inequality is exact: both terms being compared are positive and
\begin{equation}
 1877^2-64^2\,849=45625>0.
 \label{eq:spectral-certificate}
\end{equation}
Thus the proposed condition holds strictly, so the uniform vector is locally
stable under the cited second-variation test.

It is not a global maximizer.  Set
$x=(6,10,5)^{\mathsf T}$, which lies in the interior of the positive cone.
Then
\begin{equation}
 Cx=\left(\frac{142}{15},6,\frac{83}{15}\right)^{\mathsf T}.
 \label{eq:witness-image}
\end{equation}
Define
\begin{equation}
 A=\sum_i(Cx)_i^4=\frac{519655217}{15^4},
 \qquad B=\sum_i x_i^3=1341.
 \label{eq:AB}
\end{equation}
For $R=\|Cx\|_4/\|x\|_3$, one has $R^{12}=A^3/B^4$.
The conjectured MUB value is $3^{-1/12}$, whose twelfth power is $1/3$.
Their strict ordering is certified entirely in integer arithmetic:
\begin{align}
 &3(519655217)^3-15^{12}(1341)^4\notag\\
 &\hspace{2.0cm}=1409896381134395295169314>0.
 \label{eq:integer-certificate}
\end{align}
Consequently,
\begin{equation}
 \|C\|_{3\to4}\geq R>3^{-1/12}
 =3^{\,1-\lambda-\mu}.
 \label{eq:disproof}
\end{equation}
One feasible vector is sufficient for this strict lower bound; determining
the global value of $\|C\|_{3\to4}$ is unnecessary for the disproof.

\section{Consequences and limitations}

Equations~\eqref{eq:spectral-margin} and \eqref{eq:disproof} separate the
local and global questions sharply: the uniform vector passes the strict
singular-value stability test, yet another positive vector gives a larger
objective value.  Thus $\sigma_2(C)$ alone cannot certify the MUB-form norm
replacement, even for two complete rank-one qutrit measurements connected by
an exact real unitary.

This conclusion concerns only the conjectured evaluation of the induced
norm.  It does not alter the proved bound in Eq.~\eqref{eq:proved-bound}, which
remains valid with the actual norm.  Nor does the counterexample by itself
exhibit a state that violates an entropic uncertainty inequality.  It shows
instead that using the MUB value at this admissible point would require a
stronger sufficient criterion or a separate norm computation.  The example
also does not characterize the full parameter region in which the MUB formula
may remain valid.

\section*{Data and code availability}

All quantities needed for the exact argument are displayed above.  A
machine-readable certificate and a deterministic verification script using
only Python's standard library accompany the manuscript.  The script checks
orthogonality, stochasticity, spectral admissibility on the sum-zero
subspace, and the positive integer norm certificate without floating-point
arithmetic or network access.

\begin{acknowledgments}
OpenAI Codex with the gpt-5.6-sol model and Ultra reasoning setting was used
substantively for literature triage, exact-arithmetic code generation,
verification design, and manuscript drafting under prompts that fixed the
physical domain and required independent exact checks.  The reported
arithmetic was cross-checked by structurally different exact-rational
implementations.  Human author review of the science, literature, authorship
details, and wording remains pending; this local draft is not authorized for
submission.
\end{acknowledgments}

\bibliography{references}

\end{document}
